\documentclass[11pt]{article}
\usepackage[a4paper,margin=2.6cm]{geometry}
\usepackage[T1]{fontenc}
\usepackage{lmodern}
\usepackage{amsmath,amssymb,amsthm}
\newtheorem{theorem}{Theorem}
\newtheorem{lemma}[theorem]{Lemma}
\newtheorem{proposition}[theorem]{Proposition}
\theoremstyle{definition}
\newtheorem{definition}[theorem]{Definition}
\theoremstyle{remark}
\newtheorem{remark}[theorem]{Remark}
\newcommand{\dn}{{\downarrow}}
\newcommand{\cP}{\mathcal{P}}
\newcommand{\cJ}{\mathcal{J}}
\newcommand{\cC}{\mathcal{C}}
\newcommand{\cB}{\mathcal{B}}
\DeclareMathOperator{\width}{width}
\title{Counting labelled posets in time $2^{c n^2}$ with $c<1/4$\\[2pt]
\large Working note with proofs}
\author{}
\date{}
\begin{document}
\maketitle

\noindent\textbf{Status.} AI-assisted. The proofs below were written by a language
model and have not been checked by a human. Section~2 and parts of Sections~3, 6 and~8 are
machine-checked in Lean (see \texttt{lean/README.md} for the exact coverage). The counting formula of
Theorem~\ref{thm:formula} agrees with brute force for all $n\le 5$ and all $k$, in two
independently written programs. A second language-model pass refereed the proofs; its
corrections are included.

\section{Statement}

Let $\cP_n$ be the set of strict partial orders on $[n]=\{1,\dots,n\}$ and $L(n)=|\cP_n|$.
By Kleitman and Rothschild~\cite{KR75}, $\log_2 L(n) = n^2/4 + 3n/2 + O(\log n)$. Listing
all posets therefore takes time at least $2^{n^2/4}$ for large~$n$.

\begin{theorem}\label{thm:main}
$L(n)$ can be computed in time $2^{c\,n^2+O(n\log n)}$ and polynomial space, where
$c = 10-4\sqrt6 < 0.2021$.
\end{theorem}

Time is counted in bit operations; all numbers that occur have $O(n^2)$ bits.

\section{A contractible poset of wide down-sets}

Let $P$ be a finite poset. $J(P)$ is the set of down-sets of $P$, ordered by inclusion;
$\max(I)$ is the set of maximal elements of $I$; $\dn N$ is the down-set generated by~$N$.
The maps $I\mapsto\max(I)$ and $N\mapsto\dn N$ are mutually inverse bijections between
$J(P)$ and the set of antichains of~$P$. For antichains write $N\preceq N'$ if
$\dn N\subseteq\dn N'$, that is, if every element of $N$ is below or equal to some element
of~$N'$, and $N\prec N'$ if $N\preceq N'$ and $N\ne N'$.

\begin{definition}
For $k\ge1$ let $\cJ_k(P)=\{I\in J(P): |\max(I)|\ge k\}$, ordered by inclusion.
A \emph{$k$-flag} of $P$ of length $r$ is a sequence $N_0\prec N_1\prec\dots\prec N_r$ of
antichains of $P$ with $|N_i|\ge k$ for all~$i$.
\end{definition}

So $k$-flags of length $r$ correspond to chains $I_0\subsetneq\dots\subsetneq I_r$ in
$\cJ_k(P)$, and $\cJ_k(P)\ne\emptyset$ if and only if $\width(P)\ge k$.

\begin{lemma}\label{lem:submod}
For $I_1,I_2\in J(P)$,
$|\max(I_1\cup I_2)|+|\max(I_1\cap I_2)|\ \ge\ |\max I_1|+|\max I_2|$.
\end{lemma}
\begin{proof}
Let $A=\max I_1$ and $B=\max I_2$. The maximal elements of $I_1\cup I_2$ are the maximal
elements of the set $A\cup B$. Let $a\in A\cup B$ be not maximal in $A\cup B$; say $a\in A$
and $a<b$ with $b\in B$. Then $a\in I_1\cap I_2$, and $a$ is maximal in $I_1\cap I_2$
because it is maximal in~$I_1$; also $a\notin B$. An element of $A\cap B$ is maximal in
$I_1\cup I_2$ and in $I_1\cap I_2$. Hence
$|\max(I_1\cap I_2)|\ge |A\cup B|-|\max(I_1\cup I_2)|+|A\cap B|$.
\end{proof}

\begin{theorem}\label{thm:euler}
If $\width(P)\ge k$, the order complex of $\cJ_k(P)$ is contractible. Consequently
\[
  \sum_{\text{$k$-flags of $P$}} (-1)^{\text{length}}
  \;=\;\begin{cases}1,&\width(P)\ge k,\\ 0,&\text{otherwise.}\end{cases}
\]
\end{theorem}
\begin{proof}
If $\width(P)<k$ there are no $k$-flags. Otherwise let $A^*$ be an antichain of maximum
size $w\ge k$ and $I^*=\dn A^*\in\cJ_k(P)$. For $I\in\cJ_k(P)$, Lemma~\ref{lem:submod} and
$|\max(I\cap I^*)|\le w$ give $|\max(I\cup I^*)|\ge|\max I|\ge k$. So $f(I)=I\cup I^*$ is an
order-preserving map $\cJ_k(P)\to\cJ_k(P)$ with $f(I)\supseteq I$ and $f\circ f=f$, that
is, a closure operator. Hence $\cJ_k(P)$ is homotopy equivalent to $f(\cJ_k(P))$
\cite[Cor.~10.12]{B}: with $\iota$ the inclusion of the image, $\iota\circ f\ge\mathrm{id}$
and $f\circ\iota=\mathrm{id}$, and comparable order-preserving maps are
homotopic~\cite[1.3]{Q}. This image has the least element $I^*=f(I^*)$, so its order complex is a
cone. The left side of the formula is the Euler characteristic of the order complex.
\end{proof}

For $k=\width(P)$ the poset $\cJ_k(P)$ is Dilworth's distributive lattice of maximum
antichains.

\begin{remark}
The identity was also checked by computer for all $4\,473$ posets on at most $5$ labelled
points and all $k\le\width(P)$.
\end{remark}

Let $W_k(n)=|\{P\in\cP_n:\width(P)\ge k\}|$. By Theorem~\ref{thm:euler},
\begin{equation}\label{eq:W}
  W_k(n)=\sum_{P\in\cP_n}\ \sum_{\text{$k$-flags }(N_0,\dots,N_r)\text{ of }P}(-1)^r .
\end{equation}

\section{Removing the top antichain}

\begin{definition}
Let $Q$ be a poset and $A\in J(Q)$. A \emph{type} for $(Q,A)$ is a pair $(D,U)$ where
$D\subseteq A$ is a down-set of $Q$, $U\subseteq Q\setminus A$ is an up-set of $Q$, and
$d<u$ for all $d\in D$, $u\in U$. $T(Q,A)$ is the set of types.
\end{definition}

\begin{lemma}[insertion]\label{lem:ins}
Let $M\subseteq[n]$, let $Q$ be a poset on $[n]\setminus M$, $A\in J(Q)$, and
$(D_x,U_x)\in T(Q,A)$ for $x\in M$. Define a relation on $[n]$ by: the order of $Q$;
$q<x$ iff $q\in D_x$; $x<q$ iff $q\in U_x$; no relations inside~$M$. This is a strict
partial order $P$, in which $M$ is an antichain and $\dn M\setminus M=\bigcup_x D_x$.
\end{lemma}
\begin{proof}
Irreflexivity is clear, and $q<x$, $x<q$ are never both declared, because
$D_x\cap U_x\subseteq A\cap(Q\setminus A)=\emptyset$. For transitivity let $a<b<c$.
At most one of $a,b,c$ lies in $M$: adjacent elements of $M$ are unrelated, and
$a,c\in M$, $b\in Q$ would give $b\in U_a\cap D_c\subseteq(Q\setminus A)\cap A=\emptyset$.
If $a,b,c\in Q$, transitivity holds in $Q$.
If $c=x\in M$: $b\in D_x$ and $a<b$ in $Q$, so $a\in D_x$ as $D_x$ is a down-set.
If $a=x\in M$: $b\in U_x$ and $b<c$ in $Q$, so $c\in U_x$.
If $b=x\in M$: $a\in D_x$, $c\in U_x$, so $a<c$ by the definition of a type.
A transitive irreflexive relation is a strict partial order.
\end{proof}

\begin{lemma}[structure of a flag]\label{lem:struct}
Let $(N_0,\dots,N_r)$ be a $k$-flag of $P\in\cP_n$. Put $M=N_r$, $Q=P\setminus M$,
$A=\dn M\setminus M$, and for $x\in M$ let $D_x=\{q:q<x\}$, $U_x=\{q:q>x\}$. Put
$E_i=N_i\setminus M$ and $B_i=N_i\cap M$. Then
\begin{enumerate}
\item[(a)] $A\in J(Q)$, $(D_x,U_x)\in T(Q,A)$ for all $x\in M$, and $A=\bigcup_x D_x$;
\item[(b)] $B_0\subseteq B_1\subseteq\dots\subseteq B_r=M$;
\item[(c)] each $E_i$ is an antichain of $Q$ contained in $A$, and $E_r=\emptyset$;
\item[(d)] if $e\in E_i\setminus E_{i+1}$ for some $i<r$, then $e\notin E_{i'}$ for all $i'>i$.
\end{enumerate}
\end{lemma}
\begin{proof}
(a) If $q'<q\in A$ and $q<y\in M$ then $q'<y$, and $q'\notin M$ since $M$ is an antichain;
so $A$ is a down-set of $Q$. Clearly $D_x$ is a down-set, $U_x$ an up-set, $D_x\subseteq A$
and $D_x\times U_x\subseteq{<}$. If $q\in U_x\cap A$ then $x<q<y$ for some $y\in M$,
which is impossible. $A=\bigcup D_x$ by definition.
(b) Let $x\in N_i\cap M$, $i<r$. As $N_i\preceq N_{i+1}\preceq N_r$ there are
$z\in N_{i+1}$ and $y\in M$ with $x\le z\le y$. Since $M$ is an antichain, $x=z=y$.
(c) $E_i\subseteq\dn N_r\setminus M=A$.
(d) There is $z\in N_{i+1}$ with $e\le z$, and $z\ne e$, so $e<z$. For $i'>i$ there is
$y\in N_{i'}$ with $z\le y$. If $e\in N_{i'}$, then $e<y$ contradicts that $N_{i'}$ is an
antichain.
\end{proof}

By (b), with $b(x)=\min\{i:x\in N_i\}$ we have $B_i=\{x\in M:b(x)\le i\}$.
By (d), every $e\in A$ lies in $E_i$ for $i$ in an interval of $\{0,\dots,r-1\}$ (possibly
empty); we call such a sequence $(E_0,\dots,E_{r-1})$ of antichains an
\emph{interval sequence}.

\begin{definition}\label{def:data}
A \emph{configuration} for $(n,k)$ consists of: a set $M\subseteq[n]$ with $|M|=j$; a poset
$Q$ on $[n]\setminus M$; a down-set $A\in J(Q)$; an integer $r\ge0$; an interval sequence
$E=(E_0,\dots,E_{r-1})$ of antichains of $Q$ contained in $A$ (put $E_r=\emptyset$); and
for each $x\in M$ a type $(D_x,U_x)\in T(Q,A)$ and a number $b(x)\in\{0,\dots,r\}$. It is
\emph{valid} if
\begin{enumerate}
\item[(C0)] every maximal element of $A$ lies in some $D_x$;
\item[(C1)] $D_x\cap E_i=\emptyset$ whenever $b(x)\le i<r$;
\item[(C2)] for every $i<r$ and $e\in E_i$: $e\le e'$ for some $e'\in E_{i+1}$, or
            $e\in D_x$ for some $x$ with $b(x)\le i+1$;
\item[(C3)] for every $i<r$: $E_i\ne E_{i+1}$, or $b(x)=i+1$ for some $x$;
\item[(C4)] $|E_i|+|\{x:b(x)\le i\}|\ge k$ for every $i\le r$.
\end{enumerate}
\end{definition}

\begin{proposition}\label{prop:bij}
The map of Lemma~\ref{lem:struct} is a bijection from the set of pairs
$(P,\text{$k$-flag of }P)$ with $P\in\cP_n$ onto the set of valid configurations for
$(n,k)$; the flag has length~$r$.
\end{proposition}
\begin{proof}
\emph{The image is valid.} Lemma~\ref{lem:struct} shows that the image is a configuration.
(C0) holds since $A=\bigcup D_x$. (C1): $x\in N_i$ and $E_i\subseteq N_i$, and $N_i$ is an
antichain. (C2) is $N_i\preceq N_{i+1}$ for the elements of $E_i$, using
$N_{i+1}=E_{i+1}\cup\{x:b(x)\le i+1\}$. (C3) is $N_i\ne N_{i+1}$, since
$N_i\cap M\subseteq N_{i+1}\cap M$. (C4) is $|N_i|\ge k$.

\emph{The inverse.} Given a valid configuration, let $P$ be the poset of
Lemma~\ref{lem:ins} and $N_i=E_i\cup\{x:b(x)\le i\}$. $N_i$ is an antichain: $E_i$ is one
in $Q$, $M$ is one, and for $x\in N_i\cap M$, $e\in E_i$ we have $e\notin D_x$ by (C1) and
$e\notin U_x$ because $E_i\subseteq A$ and $U_x\subseteq Q\setminus A$.
$N_i\preceq N_{i+1}$: elements of $N_i\cap M$ lie in $N_{i+1}$, and elements of $E_i$ are
handled by (C2). $N_i\ne N_{i+1}$ by (C3), so $\dn N_i\subsetneq\dn N_{i+1}$.
$|N_i|\ge k$ by (C4). $N_r=M$ because $E_r=\emptyset$ and $b\le r$. So $(N_0,\dots,N_r)$
is a $k$-flag of $P$.

\emph{The two maps are mutually inverse.} Starting from $(P,\text{flag})$: $P$ is
determined by $Q$ and the sets $D_x,U_x$, and $N_i=E_i\cup B_i$. Starting from a valid
configuration: in the constructed $P$ we have $N_r=M$, $P\setminus M=Q$,
$\{q<x\}=D_x$, $\{q>x\}=U_x$, and $\dn M\setminus M=\bigcup D_x=A$ by
Lemma~\ref{lem:ins} and (C0), since a down-set contained in $A$ that contains
$\max(A)$ equals $A$. Finally $N_i\setminus M=E_i$ and $\min\{i:x\in N_i\}=b(x)$.
\end{proof}

\section{The counting formula}

Fix $Q$, $A$, $r$ and an interval sequence $E$. Let $\cC$ be the following set of
\emph{requirements}, each a pair $(e,t)$ meaning ``some $x$ has $e\in D_x$ and
$b(x)\le t$'':
\begin{itemize}
\item $(a,r)$ for each $a\in\max(A)$;
\item $(e,i+1)$ for each $i<r$ and $e\in E_i$ such that $e\not\le e'$ for all
      $e'\in E_{i+1}$.
\end{itemize}
Since $E$ is an interval sequence, each $e$ gives at most one requirement of the second
kind, so $|\cC|\le 2|A|$. Conditions (C0) and (C2) say that all requirements are met.

For $Z\subseteq\cC$ and $0\le\beta\le r$ let
\[
 c_\beta(Z)=\bigl|\{(D,U)\in T(Q,A):\ D\cap E_i=\emptyset\ (\beta\le i<r),\ \
 e\notin D\ \text{for all }(e,t)\in Z\text{ with }\beta\le t\}\bigr| .
\]
Let $\cB=\cB(E,j,k)$ be the set of vectors $(n_0,\dots,n_r)$ of non-negative integers with
sum $j$ such that $|E_i|+n_0+\dots+n_i\ge k$ for all $i\le r$, and $E_i\ne E_{i+1}$ or
$n_{i+1}\ge1$ for all $i<r$. Put
\[
 G(Q,A,E;j,k)=\sum_{Z\subseteq\cC}(-1)^{|Z|}\sum_{(n_\beta)\in\cB}
 \frac{j!}{n_0!\cdots n_r!}\prod_{\beta=0}^{r}c_\beta(Z)^{n_\beta}.
\]

\begin{lemma}\label{lem:G}
For fixed $M$ with $|M|=j$ and fixed $Q,A,r,E$, the number of valid configurations is
$G(Q,A,E;j,k)$.
\end{lemma}
\begin{proof}
By inclusion--exclusion over the set $Z$ of requirements that are \emph{not} met, the
number is $\sum_Z(-1)^{|Z|}a(Z)$, where $a(Z)$ counts the assignments
$x\mapsto(D_x,U_x,b(x))$ that satisfy (C1), (C3), (C4) and meet no requirement in $Z$.
An assignment meets no requirement in $Z$ iff for every $x$ and every $(e,t)\in Z$ with
$b(x)\le t$ we have $e\notin D_x$. Together with (C1) this is a condition on each $x$
separately, and the number of types allowed for $x$ is $c_{b(x)}(Z)$. Conditions (C3) and
(C4) depend only on the numbers $n_\beta=|\{x:b(x)=\beta\}|$ and say $(n_\beta)\in\cB$.
There are $j!/\prod n_\beta!$ functions $b$ with given $(n_\beta)$.
\end{proof}

\begin{theorem}\label{thm:formula}
For $1\le k\le n$,
\[
 W_k(n)=\sum_{j=k}^{n}\binom nj\sum_{Q\in\cP_{n-j}}\ \sum_{A\in J(Q)}\ \sum_{r=0}^{n-1}(-1)^r
 \sum_{E}G(Q,A,E;j,k),
\]
where $E$ runs over the interval sequences of length $r$ of antichains of $Q$ contained
in~$A$.
\end{theorem}
\begin{proof}
By \eqref{eq:W} and Proposition~\ref{prop:bij}, $W_k(n)$ is the sum of $(-1)^r$ over all
valid configurations. A chain in $\cJ_k(P)$ consists of distinct non-empty down-sets
ordered by inclusion, so $r\le n-1$. The number of valid configurations with given
$M,Q,A,r,E$ is given by Lemma~\ref{lem:G} and there are $\binom nj$ choices of $M$; the
increasing bijection $[n]\setminus M\to[n-j]$ identifies the posets on $[n]\setminus M$
with $\cP_{n-j}$ and preserves $G$.
\end{proof}

\section{Proof of Theorem~\ref{thm:main}}

Let $1\le k\le n$. Then $L(n)=W_k(n)+N_k(n)$, where $N_k(n)$ is the number of
$P\in\cP_n$ of width less than~$k$.

\subsection*{Wide posets}
We evaluate the formula of Theorem~\ref{thm:formula}.

\emph{Listing $\cP_m$.} Every poset on $[m]$ arises exactly once from its restriction to
$[m-1]$ by adding the point $m$ with a pair (set of points below, set of points above).
A depth-first search that tries all at most $4^{m}$ pairs at each of the $m$ levels and
keeps the transitive ones lists $\cP_m$ in time
$\sum_{i\le m}L(i-1)\,4^{i}\,m^{O(1)}\le L(m)\,2^{O(m)}$ and polynomial space, since $L$ is
non-decreasing.

\emph{Work per $Q\in\cP_m$, $m=n-j$.} There are at most $2^m$ down-sets $A$, at most $n$
values of $r$, and at most $(1+\binom{r+1}{2})^{m}\le 2^{2n\log_2 n}$ interval sequences,
since each element of $A$ occupies an interval of stages or none. There are at most
$2^{2m}$ sets $Z$ and $r+1$ values of $\beta$. Each $c_\beta(Z)$ is found by running
through the at most $4^m$ pairs $(D,U)$. The inner sum over $\cB$ is evaluated by a
dynamic programme over $\beta$ whose state is the partial sum $n_0+\dots+n_\beta\le j$;
its transitions are multiplications by $c_\beta(Z)^{n_\beta}/n_\beta!$, restricted by the
two conditions defining~$\cB$, which involve only $\beta$, the partial sum and $n_\beta$.
All of this takes time $2^{O(n\log n)}$ per~$Q$.

Hence, using~\cite{KR75}, $W_k(n)$ is computed in time
\[
\sum_{j\ge k}L(n-j)\,2^{O(n\log n)}\le L(n-k)\,2^{O(n\log n)}
=2^{(n-k)^2/4+O(n\log n)} .
\]

\subsection*{Narrow posets}
\emph{Listing.} The restriction of a poset of width less than $k$ to $[m-1]$ has width
less than~$k$. So the depth-first search above, cut off at posets of width at least $k$,
lists exactly the posets of width less than $k$. Adding a new point above all others shows
$N_k(m-1)\le N_k(m)$, so this takes time $N_k(n)\,2^{O(n)}$ and polynomial space.
Together with the wide part:

\begin{proposition}\label{prop:time}
For every $k$, $L(n)$ can be computed in time
$\bigl(L(n-k)+N_k(n)\bigr)\,2^{O(n\log n)}$ and polynomial space.
\end{proposition}

\emph{Counting.} Let $w=k-1$. An \emph{encoding} consists of
\begin{itemize}
\item an ordered list of $w'\le w$ non-empty sequences $C_1,\dots,C_{w'}$ that partition
  $[n]$ (at most $n!\,2^n$ possibilities), and
\item for every ordered pair $i\ne i'$ a non-increasing map
  $g_{ii'}\colon C_i\to\{0,\dots,|C_{i'}|\}$, where $C_i$ is ordered as a sequence.
\end{itemize}
It defines the relation in which each $C_i$ is a chain in the order of its sequence and,
for $c\in C_i$, exactly the last $g_{ii'}(c)$ elements of $C_{i'}$ are above~$c$.
The number of non-increasing maps from an $a$-chain to $\{0,\dots,b\}$ is
$\binom{a+b}{a}\le2^{a+b}$, and $\sum_{i\ne i'}(|C_i|+|C_{i'}|)=2(w'-1)n$. So there are at
most $2^{2wn+O(n\log n)}$ encodings.
Every poset of width at most $w$ has an encoding: by Dilworth's theorem it is a union of
$w'\le w$ disjoint chains, and for $c\in C_i$ the elements of $C_{i'}$ above $c$ form a
final segment of $C_{i'}$ that shrinks as $c$ increases. Hence
\begin{equation}\label{eq:narrow}
  N_k(n)\le 2^{2(k-1)n+O(n\log n)} .
\end{equation}

\subsection*{Balancing}
Take $k=\lceil\alpha n\rceil$ with $\alpha=5-2\sqrt6$, the smaller root of
$(1-\alpha)^2=8\alpha$. Then $(n-k)^2/4\le(1-\alpha)^2n^2/4=2\alpha n^2$ and
$2(k-1)n\le2\alpha n^2$, and $2\alpha=10-4\sqrt6=0.20204\ldots$ Theorem~\ref{thm:main}
follows from Proposition~\ref{prop:time}, \cite{KR75} and \eqref{eq:narrow}. \qed

\section{How far can listing the narrow posets go?}

By Proposition~\ref{prop:time}, the exponent is governed by $N_k(n)$, and \eqref{eq:narrow}
is only an upper bound. The following lower bound limits what a better count can give.

\begin{proposition}\label{prop:band}
Let $w,\ell,a\ge1$ and let $c(\ell,a)$ be the number of non-decreasing maps
$d\colon[\ell]\to\{0,\dots,\ell\}$ with $\max(0,t-2a)\le d(t)\le\max(0,t-a)$ for all~$t$.
Then $N_{w+1}(w\ell)\ge c(\ell,a)^{w(w-1)}$.
\end{proposition}
\begin{proof}
Take points $(i,t)$, $i\in[w]$, $t\in[\ell]$, and for each ordered pair $i\ne j$ such a map
$d_{ij}$. Put $(i,s)<(i,t)$ iff $s<t$, and $(j,s)<(i,t)$ iff $s\le d_{ij}(t)$ for $i\ne j$.
A relation $(j,s)<(i,t)$ implies $s<t$, and $s\le t-a$ if $i\ne j$; so the relation is
irreflexive and has no cycles. Transitivity: let $(k,u)<(j,s)<(i,t)$.
If $i,j,k$ are distinct, $u\le s-a\le t-2a$, so $t-2a\ge1$ and $u\le d_{ik}(t)$.
If $i=k\ne j$, then $u<t$. If $i=j\ne k$, then $u\le d_{ik}(s)\le d_{ik}(t)$.
If $j=k\ne i$, then $u<s\le d_{ij}(t)$. The $w$ rows are chains, so the width is at
most~$w$, and $d_{ij}(t)$ is recovered as the largest $s$ with $(j,s)<(i,t)$.
\end{proof}

For $\ell=5$, $a=1$ one has $c(5,1)=2^4$ (two choices of $d(t)$ for each $t\ge2$), so
$N_{w+1}(5w)\ge 2^{4w(w-1)}=2^{0.16\,n^2-O(n)}$ with $n=5w$, that is, for width at most
$n/5$. Since the wide part costs $2^{(1-\alpha)^2n^2/4}\ge2^{0.16n^2}$ for
$k=\alpha n\le n/5$, and $N_k$ is non-decreasing in $k$:

\begin{quote}
\emph{Any algorithm that follows Proposition~\ref{prop:time}, listing the posets of width
less than $k$, has exponent $c\ge0.16$.}
\end{quote}

So sharper counting of narrow posets can move the exponent only within
$[0.16,\,0.2021]$. Going below $0.16$ requires counting the narrow posets without listing
them.

\medskip\noindent\emph{A generalisation that fails.} In the posets of
Proposition~\ref{prop:band} there are convex subsets of about $n/3$ points that are
disjoint unions of pairwise incomparable chains, and such a subset can be inserted into
the remaining poset chain by chain. The direct analogue of Theorem~\ref{thm:euler} is
false, however: the poset of pairs $I'\subseteq I$ of down-sets for which $I\setminus I'$
is a disjoint union of pairwise incomparable chains with at least $k$ points, ordered
componentwise, has Euler characteristic $2$ for the poset $\{0,1\}<2<3$ and $k=3$
(computer search over all posets on at most $5$ points: $414$ failures in $16\,675$ cases).

\section{A sharper count of narrow posets}

Let $P(a,b)$ be the number of partial orders on the disjoint union of an $a$-chain and a
$b$-chain that extend both chains.

\begin{lemma}\label{lem:two}
$P(a,b)=\binom{a+b}{a}^2-\binom{a+b}{a-1}\binom{a+b}{a+1}
=\dfrac{(a+b)!\,(a+b+1)!}{a!\,(a+1)!\,b!\,(b+1)!}$.
\end{lemma}
\begin{proof}
Let the chains be $a_1<\dots<a_a$ and $b_1<\dots<b_b$. Such an order is given by maps
$d,e\colon[a]\to\{0,\dots,b\}$ with $b_s<a_t$ iff $s\le d(t)$ and $a_t<b_s$ iff $s>e(t)$.
Transitivity through $a_t<a_{t'}$ says that $d$ and $e$ are non-decreasing;
$b_s<a_t<b_{s'}\Rightarrow s<s'$ says $d\le e$; and $a_t<b_s<a_{t'}\Rightarrow t<t'$ then
holds automatically, since $t\ge t'$ would give $e(t)\ge e(t')\ge d(t')$. So $P(a,b)$ is the
number of pairs $d\le e$ of non-decreasing maps, that is, of pairs of non-crossing lattice
paths, and the formula is the Lindstr\"om--Gessel--Viennot determinant.
\end{proof}

\begin{lemma}[smoothing]\label{lem:smooth}
Among all $(a_1,\dots,a_w)\in\mathbb{Z}_{\ge0}^w$ with sum $n$, the product
$\prod_{i<j}P(a_i,a_j)$ is largest when $|a_i-a_j|\le1$ for all $i,j$.
\end{lemma}
\begin{proof}
From Lemma~\ref{lem:two},
\[
\frac{P(a+1,c)}{P(a,c)}=\Bigl(1+\frac{c}{a+1}\Bigr)\Bigl(1+\frac{c}{a+2}\Bigr),\qquad
\frac{P(a-1,b+1)}{P(a,b)}=\frac{a(a+1)}{(b+1)(b+2)} .
\]
The first ratio is non-increasing in $a$; the second is at least $1$ if $a\ge b+1$.
If $a_i\ge a_j+2$, replace $(a_i,a_j)$ by $(a_i-1,a_j+1)$. The factor of the pair $\{i,j\}$
does not decrease by the second identity. For $k\ne i,j$ and $c=a_k$ the two factors change
by $\frac{P(a_j+1,c)}{P(a_j,c)}\big/\frac{P(a_i,c)}{P(a_i-1,c)}\ge1$, because
$a_j\le a_i-1$. The move decreases $\sum a_i^2$, so finitely many moves lead to a balanced
profile.
\end{proof}

\begin{theorem}\label{thm:narrow}
Let $n=qw+r$ with $0\le r<w$. Then
\[
N_{w+1}(n)\le n!\,2^{n+w}\,P(q,q)^{\binom{w-r}{2}}\,P(q,q+1)^{r(w-r)}\,P(q+1,q+1)^{\binom r2}.
\]
\end{theorem}
\begin{proof}
A poset of width at most $w$ is the union of $w$ disjoint chains, some possibly empty. There
are at most $n!\binom{n+w-1}{w-1}\le n!\,2^{n+w}$ ordered $w$-tuples of disjoint sequences
covering $[n]$. For fixed chains $C_1,\dots,C_w$ the poset is determined by its restrictions
to the sets $C_i\cup C_j$, so there are at most $\prod_{i<j}P(|C_i|,|C_j|)$ posets.
Apply Lemma~\ref{lem:smooth}.
\end{proof}

With $1/\alpha=q+\theta$, $0\le\theta<1$, and $F=\log_2P$, Theorem~\ref{thm:narrow} gives
$N_{\alpha n}(n)\le2^{h_2(\alpha)n^2+O(n\log n)}$ with
\[
h_2(\alpha)=\frac{\alpha^2}{2}\Bigl[(1-\theta)^2F(q,q)+2\theta(1-\theta)F(q,q+1)
+\theta^2F(q+1,q+1)\Bigr].
\]
The equation $h_2(\alpha)=(1-\alpha)^2/4$ has the solution $\alpha=0.12260\ldots$
(with $q=8$), and Proposition~\ref{prop:time} yields:

\begin{theorem}\label{thm:main2}
$L(n)$ can be computed in time $2^{c\,n^2+O(n\log n)}$ and polynomial space with
$c<0.1925$.
\end{theorem}

Lemma~\ref{lem:two} was checked by brute force for $a+b\le5$, and
Lemma~\ref{lem:smooth} exhaustively for $n\le14$, $w\le5$.

\medskip\noindent\emph{Further.} Using exact counts of orders on three chains and
Shearer's entropy inequality, a computer-assisted argument gives $c\le0.1909$; under an
unproved log-concavity conjecture for the three-chain counts (no violation in $2101$
computed cases) it would give $c\le0.1872$. These are documented in
\texttt{narrow\_bounds/} and are not claimed here.

\medskip\noindent\emph{A better lower bound.} Proposition~\ref{prop:band} generalises: let
$\mathit{hi}\colon\{0,\dots,L\}\to\{0,\dots,L-1\}$ be non-decreasing with
$\mathit{hi}(0)=0$ and $\mathit{hi}(t)<t$ for $t\ge1$, and require
$\mathit{hi}(\mathit{hi}(t))\le d_{ij}(t)\le\mathit{hi}(t)$. The same proof applies (two
steps down reach at most $\mathit{hi}(\mathit{hi}(t))$). With
$\mathit{hi}=(0,1,1,3,3,5)$ on $t=1,\dots,6$ and chains of lengths $5$ and $6$ one gets
$18$ and $54$ admissible maps, and the bound of the previous section becomes:
\emph{any algorithm that lists the posets of width less than $k$ has exponent
$c\ge0.1651$.} So this method is confined to $0.1651\le c<0.1925$.

\section{More general identities}

The following statements are machine-checked (Lean~4 with Mathlib; see \texttt{lean/}).

\begin{theorem}[threshold]\label{thm:threshold}
Let $\mathcal L$ be a finite set of elements of a lattice, closed under joins and meets, and
let $s\colon\mathcal L\to\mathbb{Z}$ satisfy $s(p)+s(q)\le s(p\vee q)+s(p\wedge q)$. If
$\{p\in\mathcal L:s(p)\ge k\}$ is non-empty, the signed number of its non-empty chains is~$1$.
\end{theorem}
\begin{proof}
Let $m$ maximise $s$ on $\mathcal L$. Then $s(p\vee m)\ge s(p)+s(m)-s(p\wedge m)\ge s(p)$, so
$p\mapsto p\vee m$ is a closure operator on the set whose image has the least element~$m$.
\end{proof}

Theorem~\ref{thm:euler} is the case $\mathcal L=J(P)$, $s(I)=|\max I|$. Two further cases:

\begin{enumerate}
\item \emph{Pairs of down-sets.} Fix any set $R$ of pairs $x<y$. The pairs $I'\subseteq I$
of down-sets such that no $(x,y)\in R$ has $x,y\in I\setminus I'$ are closed under
componentwise union and intersection, and $s=|I\setminus I'|$ is modular. With $R$ the set
of all comparable pairs, $I\setminus I'$ is an antichain; with other choices it can be forced
to be a convex union of pairwise incomparable chains.
\item \emph{Stacks.} Tuples $I'_1\subseteq I_1\subseteq I'_2\subseteq\dots\subseteq I_h$ of
down-sets such that each $S_j=I_j\setminus I'_j$ is an antichain and, for $j<j'$, every
element of $I_j$ lies below every element outside $I'_{j'}$. They are closed under
componentwise union and intersection, and $s=\sum_j|S_j|$ is modular.
\end{enumerate}

In a stack the poset has the form $Q_0,S_1,Q_1,\dots,S_h,Q_h$ with $Q_i<Q_{i'}$ whenever
$i'\ge i+2$, and each element of $S_{j}$ is determined by a down-set of $Q_{j-1}$ and an
up-set of $Q_j$, so the layers can be inserted in product form as in
Lemma~\ref{lem:ins} (this insertion is worked out on paper only). In the posets of
Proposition~\ref{prop:band} with $\ell=5$, $a=1$, the levels $1,3,5$ form a stack with
$3n/5$ points, while the largest antichain has $n/5$.

\subsection*{What stacks would give}

Let $\sigma(P)$ be the largest number of points in the layers of a stack of $P$ (with at
most $H$ layers, $H$ fixed), and $R_k(n)$ the number of posets on $[n]$ with $\sigma<k$.
Since an antichain is a stack with one layer, $\sigma\ge\mathrm{width}$ and
$R_k(n)\le N_k(n)$.

\begin{lemma}[heredity; machine-checked]\label{lem:hered}
Let $y\in P$ have the largest strict down-set among all elements of $P$. Then every stack of
$P-y$ yields a stack of $P$ with the same layers. Hence $\sigma(P-y)\le\sigma(P)$.
\end{lemma}
\begin{proof}
$y$ is maximal. Let $D$ be its strict down-set and let $J$ be the set of indices $j$ with
$D\subseteq I'_j$ such that $I'_{j'}\supseteq P-y$ for all $j'>j$. Add $y$ to $I'_j$ and
$I_j$ for $j\in J$. The only non-trivial condition is the gap condition for $v=y$ and an
index $j'\notin J$, $j<j'$. If $D\not\subseteq I'_{j'}$, take $z\in D\setminus I'_{j'}$;
every $x\in I_j$ is below $z<y$. Otherwise some $v\ne y$ lies outside some $I'_{j''}$,
$j''>j'$. Then all of $I_{j'}\supseteq D$ is below $v$, so by the choice of $y$ the strict
down-set of $v$ equals $D$, and every $x\in I_j$, being below $v$, is below $y$.
\end{proof}

So the posets with $\sigma<k$ can be generated by one-point extensions without leaving the
class, in time $\max_{m\le n}R_k(m)\cdot2^{O(n)}$ (the statement
$\exists y\colon\sigma(P-y)\le\sigma(P)$ was also confirmed for all posets on at most $7$
points). If, in addition, the posets with $\sigma\ge k$ can be counted from the posets on
$n-k$ points as in Sections~3--4 (a sketch exists; the number of shapes of flags is
$2^{O(n\log n)}$ only for bounded $H$), then $L(n)$ is computable in time
$\bigl(L(n-k)+\max_{m\le n}R_k(m)\bigr)2^{O(n\log n)}$.

\medskip\noindent\emph{This does not yet improve Theorem~\ref{thm:main2}.} No bound on
$R_k$ is known other than $R_k\le N_k$. On the generalised band posets, where stacks formed
by whole levels can be read off, the largest entropy among families with stack fraction
below $\kappa$ equals $(1-\kappa)^2/4$ at $\kappa\approx0.24$, with value about $0.144$
(\texttt{stack\_work/floor.py}); so $0.144$ is the best this version could give on the known
examples. The family that blocks further progress consists of two blocks of pairwise
incomparable $2$-chains with free relations between the blocks; it calls for layers that
are unions of chains (case~1 above) instead of antichains.

\subsection*{Layers that are unions of chains}

Call a set of consecutive levels of a generalised band poset a \emph{window} if elements of
different rows in it are incomparable. A window is a convex union of pairwise incomparable
chains, and can be inserted chain by chain.

\begin{lemma}\label{lem:third}
Every generalised band poset on $n$ points has a stack of windows (windows
$W_1,\dots,W_m$ such that every point of level at most $\max W_j$ lies below every point of
level at least $\min W_{j+1}$) containing at least $n/3$ points.
\end{lemma}
\begin{proof}
Let $h_0=L$ and $h_{i+1}=\mathit{hi}(h_i)$, and let $B_i$ be the set of levels in
$(h_i,h_{i-1}]$. $B_i$ is a window, because $\mathit{hi}(t)\le h_i$ for $t\le h_{i-1}$.
For $s=h_i+1=\min B_i$ we have
$\mathit{hi}(\mathit{hi}(s))\ge\mathit{hi}(\mathit{hi}(h_i))=h_{i+2}=\max B_{i+3}$,
so every point of level at most $\max B_{i+3}$ lies below every point of level at least
$\min B_i$. Hence the blocks $B_i$ with $i$ in a fixed residue class modulo $3$ form a stack
of windows, and one of the three classes contains at least a third of the levels.
\end{proof}

So for all the constructions known to us, a method based on stacks of chain-union layers
would need only the posets on $2n/3$ points. What is missing is the same as before: a bound
on the number of posets without such a stack.

\subsection*{The frontier by rank}

The transfer matrix of our earlier paper on distributive lattices also counts posets by
number of points. Build $P$ in the order of $|{\downarrow}x|$. Let
$X_r=\{x:|{\downarrow}x|<r\}$, $b=|X_r|$, and call $z\in X_r$ \emph{dead} if
$|{\uparrow}z\cap X_r|\ge b-r+2$.

\begin{lemma}
A dead element lies below every element outside $X_r$. If $z$ is not dead and $x\in X_r$,
then $|{\uparrow}z\cap X_r|+|{\downarrow}x|\le b$.
\end{lemma}
\begin{proof}
Let $y\notin X_r$ have the smallest $|{\downarrow}y|$ among the elements outside $X_r$ that
are not above $z$. Its strict down-set $D$ has at least $r-1$ elements and avoids
${\uparrow}z$; by the choice of $y$ it lies in $X_r\setminus{\uparrow}z$, which has at most
$r-2$ elements. The second statement is the definition, with $|{\downarrow}x|\le r-1$.
\end{proof}

So a state only needs the poset induced on the live elements (and counters), and that poset
satisfies $\max|{\uparrow}z|+\max|{\downarrow}x|\le n$. Counting the up-staircases of
$w$ chains under this constraint with a generating-function bound gives at most
$2^{f(\alpha)n^2}$ live posets of width below $\alpha n$, with
$f(\alpha)=\min_{0<t<1/2}\bigl[\alpha\log_2\frac1t
+\alpha^2\log_2\bigl(1+\frac{t^2}{1-2t}\bigr)\bigr]$, and
$f(\alpha)=(1-\alpha)^2/4$ at $\alpha\approx0.136$ with value $0.1865$. This would improve
Theorem~\ref{thm:main2} slightly, but the transfer matrix for labelled posets with a width
bound has not been written out or tested; only the lemma was checked, on all posets with at
most $7$ points.

\section{Narrow posets and templates}

This section reduces the count of narrow posets to a finite extremal problem. It is a
proof sketch that has not yet been refereed; the numerical evaluation is in progress.

\subsection*{Templates}
Let $C_1,\dots,C_m$ be disjoint chains of lengths $a_1,\dots,a_m$. An order extending them
is given by non-decreasing maps $d_{ij}\colon[a_i]\to\{0,\dots,a_j\}$ for $i\ne j$: the
element at position $t$ of $C_i$ is above exactly the first $d_{ij}(t)$ elements of $C_j$
(put $d(0)=0$). Transitivity through a third chain reads
\begin{equation}\label{eq:tri}
  d_{ik}(t)\ \ge\ d_{jk}(d_{ij}(t))\qquad(i,j,k\text{ distinct}).
\end{equation}
A \emph{template} assigns to each pair $\{i,j\}$ a non-empty set $A_{ij}$ of orders on
$C_i\cup C_j$ such that every choice of one member from each set gives an order on
$C_1\cup\dots\cup C_m$. Its \emph{entropy} is $\sum_{i<j}\log_2|A_{ij}|$, and
$\tau_m(a_1,\dots,a_m)$ is the largest entropy of a template. The band posets of
Proposition~\ref{prop:band} are the realisations of templates.

\begin{lemma}[envelopes]\label{lem:env}
Let $H_{ij}$ be the pointwise maximum of $d_{ij}$ over $A_{ij}$. Then every member of
$A_{ik}$ satisfies $d_{ik}\ge H_{jk}\circ H_{ij}$ for every third chain~$j$. Hence
$\tau_3(a_1,a_2,a_3)\le\tau_3^{\rm up}(a_1,a_2,a_3)$, the maximum over six non-decreasing
maps $H_{ij}$ of $\sum_{\{i,k\}}\log_2$ of the number of orders $(d_{ik},d_{ki})$ on
$C_i\cup C_k$ with $H_{jk}\circ H_{ij}\le d_{ik}\le H_{ik}$ and
$H_{ji}\circ H_{kj}\le d_{ki}\le H_{ki}$.
\end{lemma}
\begin{proof}
Fix $t$. Choose a member of $A_{ij}$ with $d_{ij}(t)=H_{ij}(t)=:s$ and a member of $A_{jk}$
with $d_{jk}(s)=H_{jk}(s)$. Together with any member of $A_{ik}$ they form an order, so
\eqref{eq:tri} gives $d_{ik}(t)\ge H_{jk}(H_{ij}(t))$.
\end{proof}

\subsection*{The counting theorem}
We use the following result of Falgas-Ravry, O'Connell and Uzzell~\cite{FOU}. A property
$\mathcal Q=(\mathcal Q_B)_B$ of colourings of the edges of $K_B$ (vertex set $[B]$) with a
fixed finite set of colours is \emph{order-hereditary} if it is closed under restriction to
subsets of $[B]$ with the induced order. A template for $K_B$ gives each edge a non-empty
palette; $\mathrm{ex}(B,\mathcal Q)$ is the largest entropy (here in bits) of a template all of
whose realisations lie in $\mathcal Q_B$. Then $\mathrm{ex}(B,\mathcal Q)/\binom B2$ is
non-increasing with limit $\pi(\mathcal Q)$ \cite[Prop.~2.10]{FOU}, and for every
$\varepsilon>0$ and all large $B$,
$\log_2|\mathcal Q_B|\le(\pi(\mathcal Q)+\varepsilon)\binom B2$ \cite[Cor.~2.15]{FOU}.
In particular $\pi(\mathcal Q)\le\mathrm{ex}(3,\mathcal Q)/3$.

\begin{theorem}\label{thm:templates}
For all $K\ge1$ and $\varepsilon>0$ there is $w_0$ such that for all $w\ge w_0$ and all
disjoint chains $C_1,\dots,C_w$ of lengths $a_i\le K$, the number of orders extending them is
at most
\[
 2^{\,\frac{w^2}{6}\,\mathbb{E}\,\tau_3(X,Y,Z)+\varepsilon w^2},
\]
where $X,Y,Z$ are independent with the empirical distribution of $(a_1,\dots,a_w)$.
\end{theorem}
\begin{proof}[Proof sketch]
\emph{Bundles.} Fix a large $g$. Let $w_\ell$ be the number of chains of length $\ell$,
$\gamma_\ell=\lfloor gw_\ell/w\rfloor$, $g'=\sum_\ell\gamma_\ell\in[g-K-1,g]$ and
$B=\lfloor w/g\rfloor$. Form $B$ \emph{bundles}, each consisting of $\gamma_\ell$ chains of
length $\ell$ for every $\ell$. At most $w(K+1)/g+g$ chains are left over; their relations
to all chains contribute at most $4K\,w\,(w(K+1)/g+g)$ bits, which is at most
$\varepsilon w^2/3$ for $g$ large and then $w$ large.

\emph{The property.} Colour the edge $ij$ ($i<j$) of $K_B$ by the order induced on the union
of bundles $i$ and $j$; the colour set is finite and depends only on $(\gamma_\ell)$.
Let $\mathcal Q_B$ consist of the colourings that come from an order on the union of all
bundles. This is order-hereditary. There are finitely many vectors $(\gamma_\ell)$ for given
$g,K$, so the constant in \cite[Cor.~2.15]{FOU} can be chosen uniformly.

\emph{Three bundles.} Let $t$ be a template on $K_3$ with all realisations in $\mathcal Q_3$.
Each vertex lies in two edges whose colours are chosen independently, so the order inside
each bundle is the same for all colours of all palettes. A colour of the edge $ij$ is
therefore determined by the orders on the pairs $C_x\cup C_y$, $x$ in bundle~$i$, $y$ in
bundle~$j$, and $|t(ij)|\le\prod_{x,y}|A_{xy}|$, where $A_{xy}$ is the set of orders on
$C_x\cup C_y$ occurring in $t(ij)$. For $x,y,z$ in three different bundles every
combination from $A_{xy}\times A_{xz}\times A_{yz}$ occurs in a realisation, so
$(A_{xy},A_{xz},A_{yz})$ is a template and
$\log_2|A_{xy}|+\log_2|A_{xz}|+\log_2|A_{yz}|\le\tau_3(a_x,a_y,a_z)$.
Each pair $(x,y)$ lies in $g'$ such triples, so
$\mathrm{ex}(3,\mathcal Q)\le\frac1{g'}\sum_{x,y,z}\tau_3(a_x,a_y,a_z)
=g'^2\,\mathbb{E}\,\tau_3(X',Y',Z')$ with $X',Y',Z'$ distributed as the lengths in a bundle.

\emph{Conclusion.} $\log_2|\mathcal Q_B|\le\binom B2\bigl(\tfrac{g'^2}{3}\mathbb{E}\,\tau_3(X',Y',Z')
+\varepsilon'\bigr)$, $Bg'\le w$, and the distribution of the lengths in a bundle differs
from the empirical one by at most $(K+1)/g$ in each coordinate, while
$\tau_3\le 3\log_2P(K,K)\le12K$.
\end{proof}

The same argument with $m$ bundles gives $\frac{w^2}{2}\,\mathbb{E}\,\tau_m/\binom m2$, and
these bounds decrease to the truth as $m\to\infty$, because
$\mathrm{ex}(m,\mathcal Q)/\binom m2\to\pi(\mathcal Q)$ and
$|\mathcal Q_B|\ge2^{\pi(\mathcal Q)\binom B2}$.

\medskip\noindent\emph{Consequence.} After cutting long chains as in the previous section,
\[
h(\alpha)\ \le\ \frac{\tilde\alpha^{2}}{6}\,\sup\bigl\{\mathbb{E}_{\mu^{3}}\tau_3:\ 
\mu\text{ on }\{0,\dots,K\}\text{ with mean }1/\tilde\alpha\bigr\},\qquad
\tilde\alpha=\alpha+1/K,
\]
where $2^{h(\alpha)n^2}$ is the number of posets of width at most $\alpha n$. The
three-chain \emph{count} used before is replaced by the three-chain \emph{template}, which
is much smaller: a randomised search (a lower estimate of $\tau_3$, not a bound) gives
$8.73$ bits per pair for three chains of length $5$, against $14.24$ for the pair count,
about $12.7$ for the three-chain count and $8.34$ for the band construction.

\section{Remarks}

\begin{enumerate}
\item The constant is not optimised. The bound for narrow posets is crude, and the full
order complex could perhaps be replaced by a smaller complex.
\item The result is about labelled posets. For unlabelled posets one would need an
equivariant form of Theorem~\ref{thm:euler}; this has not been worked out.
\item The $2^{O(n\log n)}$ factor makes the method useless for the values of $n$ within
reach of computation ($n\approx20$).
\item Whether Theorem~\ref{thm:euler} is known has not been checked.
\item Numerical check of Theorem~\ref{thm:formula}: the script
\texttt{test\_flag\_formula.py} evaluates the formula and compares with brute force.
\end{enumerate}

\begin{thebibliography}{9}
\bibitem{B} A.~Bj\"orner, Topological methods, in: \emph{Handbook of Combinatorics},
  Vol.~2, Elsevier, 1995, 1819--1872.
\bibitem{FOU} V.~Falgas-Ravry, K.~O'Connell and A.~Uzzell, Multicolour containers, extremal
  entropy and counting, \emph{Random Structures Algorithms} \textbf{54} (2019), 676--720;
  arXiv:1801.05195.
\bibitem{KR75} D.~J.~Kleitman and B.~L.~Rothschild, Asymptotic enumeration of partial
  orders on a finite set, \emph{Trans.\ Amer.\ Math.\ Soc.} \textbf{205} (1975), 205--220.
\bibitem{Q} D.~Quillen, Homotopy properties of the poset of nontrivial $p$-subgroups of a
  group, \emph{Adv.\ Math.} \textbf{28} (1978), 101--128.
\end{thebibliography}
\end{document}
